\section{Game domain and rulesets}
\label{sec:game}

\subsection{CellGameModule}
Primary gameplay module. Owns a \symbolref{CellContext} and an enum
\symbolref{CellState}:
\begin{center}
\texttt{NORMAL} $\mid$ \texttt{EDIT} $\mid$ \texttt{SAVE} $\mid$ \texttt{LOAD} $\mid$ \texttt{EXIT}
\end{center}

\begin{statusbox}
Mode handling is an \textbf{enum + switch} inside the module, not the older
class-per-state hierarchy (archived under \pathref{archive/old-states/}).
SAVE/LOAD cases may still be incomplete relative to method helpers.
\end{statusbox}

\begin{inferencebox}
Earlier design used State classes so each mode had its own ``main loop'' body.
The enum approach is simpler for a small mode set; reintroduce State objects
only if modes grow heavy (UI stacks, async load, etc.).
\end{inferencebox}

\subsection{CellContext}
Constructs \symbolref{Canvas} from env vars (\texttt{CanvasX}, \texttt{CanvasY})
and selects a \symbolref{RuleSet} from a mode string
(\texttt{GAME\_OF\_LIFE}, \texttt{BRIANS\_BRAIN}, \ldots).

\subsection{Canvas}
\begin{itemize}
  \item CPU grid: \symbolref{lifeCanvas} (and RGB staging \symbolref{texCanvasBuffer}).
  \item Currently also a \symbolref{Drawable} with GL texture id --- migration
        target is to split this (\cref{sec:rendering-target}).
\end{itemize}

\subsection{Rulesets}
Classic polymorphic hierarchy:
\begin{lstlisting}
class RuleSet {
  virtual void evalCell(...) const;
  virtual void evaluateNeighbors(...) const;
  void calcGeneration(...); // shared sweeping
};
\end{lstlisting}
Concrete rules live in \pathref{Rulesets/}. This package is relatively clean:
depends on canvas/domain, not on windowing.

\begin{decisionbox}[title={Keep rules pure}]
Rulesets should stay free of rendering and input. They read/write cell buffers
only. That keeps them testable and migration-safe.
\end{decisionbox}
